\section{Utility Maximization}

\subsection{Basics for general $\Omega$}

Let me first explain all the objects in normal language.

\begin{itemize}
    \item $x>0$ is the initial capital, i.e. the money we start with.
    \item $\theta$ is the trading strategy. It tells us how many units of the risky assets we hold over time.
    \item $G(\theta)$ is the gains process generated by the strategy $\theta$. So $G_t(\theta)$ is the profit or loss from trading up to time $t$.
    \item $V(x,\theta)$ is the wealth process. It tracks my total wealth over time when I start with $x$ and trade according to $\theta$.
    \item Self-financing means that after time $0$ we do not add money to the portfolio and we do not take money out. All changes in wealth come only from trading gains and losses.
    \item $0$-admissible means that the wealth is never allowed to become negative.
\end{itemize}

So morally:
\[
    V_t(x,\theta)=x+G_t(\theta).
\]
For example, if $x=100$ and after trading we have gained $20$ by time $T$, then
\[
    V_T=100+20=120.
\]
If instead we lose $30$, then
\[
    V_T=100-30=70.
\]
The condition of $0$-admissibility says that values like $V_t=-5$ are not allowed.

The utility function $U$ measures how happy we are with terminal wealth. It is constructed so that:
\begin{itemize}
    \item more money is always better;
    \item the additional happiness from extra money becomes smaller when we are already rich;
    \item getting an extra $100$ is more important if we currently have $100$ than if we already have one million;
    \item near $0$, money is extremely important;
    \item near infinity, one additional unit of money does not matter much anymore.
\end{itemize}

The small letter $u(x)$ represents the best possible expected utility when we start with initial capital $x$. In other words, among all admissible trading strategies, we choose the one which maximizes
\[
    \E[U(\text{final wealth})].
\]

The set $\C(x)$ tells us that it is enough to look at possible terminal payoffs. We do not always need to look at the full wealth process at intermediate times. We can reformulate the problem as maximizing utility over all possible final payoffs.

Finally, the Legendre transform helps us pass from the primal problem to a dual problem. The primal problem is a maximization problem. The dual problem is a minimization problem, and it is often easier to solve.

\subsubsection*{Wealth processes}

Since we restrict to utility functions on $\R_+=(0,\infty)$, we consider nonnegative wealth processes.

\begin{definicija}
For $x>0$, define
\[
    \V(x)
    :=
    \left\{
        V=V(x,\theta): V_t=x+G_t(\theta),\ t\in[0,T],
        \text{ where } \theta \text{ is self-financing and } V_t\geq 0
    \right\}.
\]
Equivalently,
\[
    \V(x)
    =
    \left\{
        x+G(\theta): \theta \text{ self-financing and } x+G_t(\theta)\geq 0
        \text{ for all } t
    \right\}.
\]
The condition
\[
    x+G_t(\theta)\geq 0
\]
is called $0$-admissibility.
\end{definicija}

\begin{primer}
Suppose $x=100$ and at terminal time $T$ a strategy gives gain
\[
    G_T(\theta)=25.
\]
Then the final wealth is
\[
    V_T=100+25=125.
\]
This is admissible if the whole wealth path stays nonnegative, not just the final value.

If at some intermediate time $t$ we had
\[
    V_t=100+G_t(\theta)=-10,
\]
then this strategy would not be $0$-admissible.
\end{primer}

\subsubsection*{Utility functions}

\begin{definicija}
A \textit{utility function} is a function
\[
    U:\R_+ \to \R
\]
such that:
\begin{itemize}
    \item $U$ is strictly increasing;
    \item $U$ is strictly concave;
    \item $U\in C^1$;
    \item $U$ satisfies the \textit{Inada conditions}
    \[
        U'(0):=\lim_{x\downarrow 0} U'(x)=+\infty,
    \]
    and
    \[
        U'(\infty):=\lim_{x\to\infty} U'(x)=0.
    \]
\end{itemize}
\end{definicija}

\begin{primer}
The function
\[
    U(x)=\log x
\]
is a classical example. It is increasing and concave. Also,
\[
    U'(x)=\frac{1}{x},
\]
so
\[
    U'(0)=+\infty,
    \qquad
    U'(\infty)=0.
\]
This matches the intuition: if we have almost no money, a little more money is very valuable; if we are extremely rich, one more unit matters almost not at all.
\end{primer}

\begin{primer}
Another classical example is power utility:
\[
    U(x)=\frac{1}{\gamma}x^\gamma,
    \qquad 0<\gamma<1.
\]
Then
\[
    U'(x)=x^{\gamma-1}.
\]
Since $\gamma-1<0$, we get
\[
    \lim_{x\downarrow 0} x^{\gamma-1}=+\infty,
    \qquad
    \lim_{x\to\infty} x^{\gamma-1}=0.
\]
So power utility also satisfies the Inada conditions.
\end{primer}

\begin{opomba}
By monotonicity, the limit
\[
    U(0):=\lim_{x\downarrow 0}U(x)
\]
always exists in $[-\infty,\infty)$. Therefore, for $V\in\V(x)$, the random variable $U(V_T)$ is well-defined in $[-\infty,\infty]$ almost surely.

We define
\[
    \E[U(V_T)] := -\infty
\]
if the negative part of $U(V_T)$ is not integrable, i.e. if
\[
    U(V_T)^- \notin L^1(\Pp).
\]
This convention prevents undefined expressions of the form $\infty-\infty$.
\end{opomba}

\subsubsection*{The primal utility maximization problem}

Our problem is to find the indirect utility function
\[
    u(x)
    :=
    \sup_{V\in\V(x)}
    \E[U(V_T)].
\]
So $u(x)$ is the best expected utility we can obtain from initial capital $x$.

Since we can always choose the strategy $\theta=0$, meaning ``do nothing'', we have
\[
    V_t=x
    \quad\text{for all }t,
\]
and therefore
\[
    u(x)\geq U(x).
\]

In general,
\[
    u(x)\in(-\infty,\infty].
\]

\subsubsection*{Standing assumption}

We assume that
\[
    \exists x_0>0
    \quad\text{such that}\quad
    u(x_0)<\infty.
\]
This assumption excludes degenerate cases where the investor can obtain infinite expected utility.

\begin{primer}
If this assumption holds, then under the usual convexity assumptions on the set of attainable wealths, the map
\[
    x\mapsto u(x)
\]
is increasing and concave. Moreover, finiteness at one point often implies finiteness for all $x>0$ in the standard utility maximization setting.
\end{primer}

\begin{primer}
If $U$ is unbounded from above and the market allows arbitrage, then one typically gets
\[
    u(x)\equiv \infty.
\]
The reason is that arbitrage allows us to improve wealth without risk, and if utility is unbounded, we can push expected utility arbitrarily high.
\end{primer}

\subsection{Legendre transform}

The dual problem is built using the convex conjugate of the utility function.

\begin{definicija}
For $y>0$, define
\[
    J(y)
    :=
    \sup_{x>0}
    \bigl(U(x)-xy\bigr).
\]
This is the Legendre-Fenchel transform of $U$, with the sign convention used in utility maximization.

The variable $y$ can be interpreted as a dual variable, or morally as a ``price of wealth''.
\end{definicija}

\begin{primer}
For a fixed $y>0$, the expression
\[
    U(x)-xy
\]
has two parts:
\begin{itemize}
    \item $U(x)$ is the happiness from wealth $x$;
    \item $xy$ is a linear penalty for choosing wealth $x$ when the dual price is $y$.
\end{itemize}
The Legendre transform $J(y)$ chooses the value of $x$ which gives the best tradeoff between utility and penalty.
\end{primer}

\begin{lema}
Let $U$ be a utility function. Then
\[
    J:\R_+\to\R
\]
is strictly decreasing, strictly convex, and belongs to $C^1$.

Moreover,
\[
    J'(y)=-(U')^{-1}(y).
\]
If we define
\[
    I(y):=(U')^{-1}(y),
\]
then
\[
    J'(y)=-I(y)
\]
and
\[
    J(y)=U(I(y))-yI(y).
\]

The conjugacy relation is
\[
    U(x)
    =
    \inf_{y>0}
    \bigl(J(y)+xy\bigr),
    \qquad x>0.
\]

Also, in the extended sense,
\[
    J(0)=U(\infty):=\lim_{x\to\infty}U(x),
\]
and
\[
    J(\infty)=U(0):=\lim_{x\downarrow 0}U(x).
\]
\end{lema}

\begin{proof}
For fixed $x>0$, the function
\[
    y\mapsto U(x)-xy
\]
is affine and decreasing in $y$. The supremum of affine functions is convex, which gives convexity of $J$. The strict monotonicity and differentiability follow from the strict concavity and differentiability of $U$.

To find the maximizer in
\[
    J(y)=\sup_{x>0}(U(x)-xy),
\]
differentiate with respect to $x$:
\[
    U'(x)-y=0.
\]
Hence the optimizer is
\[
    x=(U')^{-1}(y)=I(y).
\]
Substituting this back gives
\[
    J(y)=U(I(y))-yI(y).
\]
By the envelope theorem, or by differentiating this expression, we get
\[
    J'(y)=-I(y).
\]
The remaining statements are standard facts from convex analysis.
\end{proof}

\begin{primer}
Let
\[
    U(x)=\log x.
\]
Then
\[
    U'(x)=\frac{1}{x}.
\]
Solving
\[
    U'(x)=y
\]
gives
\[
    \frac{1}{x}=y,
    \qquad
    x=\frac{1}{y}.
\]
Therefore
\[
    I(y)=\frac{1}{y}.
\]
Now
\[
    J(y)
    =
    U(I(y))-yI(y)
    =
    \log\left(\frac{1}{y}\right)-y\frac{1}{y}
    =
    -\log y-1.
\]
So for logarithmic utility,
\[
    J(y)=-\log y-1.
\]
\end{primer}

\begin{primer}
Let
\[
    U(x)=\frac{1}{\gamma}x^\gamma,
    \qquad 0<\gamma<1.
\]
Then
\[
    U'(x)=x^{\gamma-1}.
\]
Solving
\[
    x^{\gamma-1}=y
\]
gives
\[
    I(y)=y^{\frac{1}{\gamma-1}}.
\]
Thus
\[
    J(y)
    =
    U(I(y))-yI(y)
    =
    \frac{1}{\gamma}y^{\frac{\gamma}{\gamma-1}}
    -
    y^{1+\frac{1}{\gamma-1}}.
\]
Since
\[
    1+\frac{1}{\gamma-1}
    =
    \frac{\gamma}{\gamma-1},
\]
we get
\[
    J(y)
    =
    \left(\frac{1}{\gamma}-1\right)y^{\frac{\gamma}{\gamma-1}}
    =
    \frac{1-\gamma}{\gamma}
    y^{\frac{\gamma}{\gamma-1}}.
\]
So
\[
    J(y)=\frac{1-\gamma}{\gamma}y^{\frac{\gamma}{\gamma-1}}.
\]
\end{primer}

\begin{opomba}
The logarithmic utility case can be seen as a limiting case of power utility.
\end{opomba}

\begin{naloga}
Some useful exercises and facts:
\begin{enumerate}
    \item If $U$ is increasing and concave, but not strictly concave, then $u$ is still increasing and concave. Strict concavity of $u$ may fail.
    \item If $U$ is strictly increasing and the market satisfies no-arbitrage assumptions, then we do not expect utility to explode due to free profit.
    \item If the market admits arbitrage and $U$ is unbounded from above, then usually
    \[
        u(x)=\infty.
    \]
\end{enumerate}
\end{naloga}

\subsection{Abstract formulation of the primal problem and the dual problem}

We want to solve the primal problem
\[
    u(x)=\sup_{V\in\V(x)}\E[U(V_T)].
\]
Here $\V(x)$ denotes all nonnegative wealth processes starting from initial capital $x$.

Instead of optimizing over whole wealth processes, we can optimize over their possible terminal values.

\begin{definicija}
Define
\[
    \C(x)
    :=
    \bigl(x+G_T(\Theta_{\mathrm{adm}})-L_+^0\bigr)
    \cap L_+^0.
\]
Equivalently,
\[
    \C(x)
    =
    \left\{
        f\in L_+^0:
        f=x+G_T(\theta)-Y,
        \ \theta\in\Theta_{\mathrm{adm}},
        \ Y\in L_+^0
    \right\}.
\]
So $\C(x)$ is the set of all nonnegative terminal payoffs which can be superreplicated from initial capital $x$.

In other words, $f\in\C(x)$ means:
there exists an admissible strategy $\theta$ such that
\[
    f\leq x+G_T(\theta).
\]
\end{definicija}

\begin{primer}
Suppose $x=100$ and some admissible strategy gives terminal wealth
\[
    V_T=130.
\]
Then every nonnegative payoff $f$ satisfying
\[
    0\leq f\leq 130
\]
belongs to $\C(100)$, because the strategy produces enough money to cover $f$.

For example,
\[
    f=120
\]
is in $\C(100)$, and so is
\[
    f=80.
\]
The idea is that if we can end with $130$, then we can also superreplicate a smaller payoff.
\end{primer}

Then we have the equivalent formulation
\[
    u(x)
    =
    \sup_{f\in\C(x)}
    \E[U(f)].
\]

Indeed, the set of terminal wealths
\[
    \V_T(x):=\{V_T:V\in\V(x)\}
\]
is contained in $\C(x)$, because every terminal wealth can trivially superreplicate itself.

Conversely, if
\[
    f\in\C(x),
\]
then there exists $V\in\V(x)$ such that
\[
    f\leq V_T.
\]
Since $U$ is increasing,
\[
    U(f)\leq U(V_T).
\]
Therefore
\[
    \E[U(f)]
    \leq
    \E[U(V_T)]
    \leq
    u(x).
\]
Taking the supremum over $f\in\C(x)$ gives
\[
    \sup_{f\in\C(x)}\E[U(f)]
    \leq u(x).
\]
The reverse inequality is immediate because terminal wealths are included in $\C(x)$. Hence
\[
    u(x)
    =
    \sup_{f\in\C(x)}
    \E[U(f)].
\]

\begin{opomba}
This is useful because the problem can now be viewed as an optimization problem over terminal payoffs, rather than over entire stochastic processes.
\end{opomba}

\subsection{Motivation for the dual problem}

Our goal is to solve the primal problem
\[
    u(x)=\sup_{V\in\V(x)}\E[U(V_T)].
\]
We will do this using a dual problem involving martingale deflators, or in special cases equivalent local martingale measures.

Assume that the set of equivalent local martingale measures is nonempty. Let
\[
    \Q\sim\Pp
\]
be such a measure, and denote its density process by
\[
    Z_t^{\Q,\Pp}
    :=
    \E_{\Pp}\left[\frac{d\Q}{d\Pp}\,\middle|\,\F_t\right].
\]

Take any wealth process
\[
    V=x+G(\theta)\in\V(x).
\]
Under $\Q$, the gains process is a local martingale. Since $V$ is nonnegative, it is a $\Q$-supermartingale. By Bayes' theorem, the product
\[
    Z^{\Q,\Pp}V
\]
is a $\Pp$-supermartingale.

This motivates the following definition.

\begin{definicija}
For $z>0$, define
\[
    \Z(z)
    :=
    \left\{
        Z:
        Z \text{ is nonnegative, adapted, } Z_0=z,
        \text{ and } ZV \text{ is a } \Pp\text{-supermartingale for all } V\in\V(1)
    \right\}.
\]
The elements of $\Z(z)$ are called supermartingale deflators.
\end{definicija}

\begin{primer}
If $\F_0$ is trivial and $\Q$ is an equivalent local martingale measure, then the density process
\[
    Z_t^{\Q,\Pp}
    =
    \E_{\Pp}\left[\frac{d\Q}{d\Pp}\,\middle|\,\F_t\right]
\]
belongs to $\Z(1)$.

So $\Z(1)$ contains the usual martingale density processes, but it may also contain more general objects.
\end{primer}

\begin{opomba}
Since $\V(1)$ contains the constant wealth process
\[
    V_t\equiv 1,
\]
every $Z\in\Z(z)$ is itself a nonnegative $\Pp$-supermartingale.

From now on, we assume for simplicity that $\F_0$ is trivial. Then $Z_0=z$ is just a deterministic constant.
\end{opomba}

Now take
\[
    V\in\V(x),
    \qquad
    Z\in\Z(z).
\]
Since $V/x\in\V(1)$, the product
\[
    Z\frac{V}{x}
\]
is a supermartingale. Equivalently, $ZV$ is a supermartingale starting at
\[
    Z_0V_0=zx.
\]
Therefore
\[
    \E[Z_TV_T]\leq xz.
\]

By definition of $J$, for all $a>0$ and $y>0$,
\[
    J(y)\geq U(a)-ay.
\]
Rearranging,
\[
    U(a)\leq J(y)+ay.
\]
Taking
\[
    a=V_T,
    \qquad
    y=Z_T,
\]
we get
\[
    U(V_T)
    \leq
    J(Z_T)+Z_TV_T.
\]
Therefore
\[
    \E[U(V_T)]
    \leq
    \E[J(Z_T)]
    +
    \E[Z_TV_T].
\]
Using the supermartingale inequality,
\[
    \E[Z_TV_T]\leq xz,
\]
we obtain
\[
    \E[U(V_T)]
    \leq
    \E[J(Z_T)]
    +
    xz.
\]
Now take the supremum over all $V\in\V(x)$:
\[
    u(x)
    \leq
    \E[J(Z_T)]
    +
    xz.
\]
Finally, take the infimum over $Z\in\Z(z)$.

\begin{definicija}
Define the dual value function
\[
    j(z)
    :=
    \inf_{Z\in\Z(z)}
    \E[J(Z_T)].
\]
Then
\[
    u(x)
    \leq
    j(z)+xz,
    \qquad x,z>0.
\]
\end{definicija}

\begin{opomba}
The primal problem maximizes a concave functional:
\[
    \sup_{V\in\V(x)}\E[U(V_T)].
\]
The dual problem minimizes a convex functional:
\[
    \inf_{Z\in\Z(z)}\E[J(Z_T)].
\]
This is the central primal-dual picture.
\end{opomba}

\subsection{Abstract dual formulation}

To formulate the dual problem only in terms of terminal random variables, define
\[
    \D(z)
    :=
    \left\{
        h\in L_+^0:
        h\leq Z_T
        \text{ for some } Z\in\Z(z)
    \right\}.
\]

\textbf{\color{red}Important:} This is the dual analogue of the set $\C(x)$.

The set $\C(x)$ contains terminal payoffs which are dominated by admissible terminal wealths:
\[
    f\leq V_T.
\]
The set $\D(z)$ contains terminal dual variables which are dominated by terminal values of deflators:
\[
    h\leq Z_T.
\]

Because $J$ is decreasing, if
\[
    h\leq Z_T,
\]
then
\[
    J(h)\geq J(Z_T).
\]
Therefore one has to be careful with the direction of inequalities. In the standard abstract formulation, $\D(z)$ is chosen so that the dual minimization problem becomes
\[
    j(z)
    =
    \inf_{h\in\D(z)}
    \E[J(h)].
\]
The precise definition of $\D(z)$ may vary slightly depending on the convention used in the lecture notes. The important point is that $\D(z)$ is the terminal-value version of the dual domain.

So the abstract primal-dual formulation is:

\[
    u(x)
    =
    \sup_{f\in\C(x)}
    \E[U(f)],
\]
and
\[
    j(z)
    =
    \inf_{h\in\D(z)}
    \E[J(h)].
\]

\begin{primer}
Think of the two sets like this:
\begin{itemize}
    \item $\C(x)$ contains things we can afford from initial wealth $x$.
    \item $\D(z)$ contains possible dual prices or state-price densities starting from $z$.
\end{itemize}

The primal problem asks:
\[
    \text{What is the best payoff I can afford?}
\]
The dual problem asks:
\[
    \text{What is the cheapest dual object that supports this optimization?}
\]
\end{primer}

\subsection{Solving the dual problem}

The goal is to show that the abstract dual problem
\[
    j(z)
    =
    \inf_{h\in\D(z)}
    \E[J(h)]
\]
has a unique solution
\[
    h_z^*\in\D(z)
\]
whenever
\[
    j(z)<\infty.
\]

The basic intuition is the following.

If we were minimizing a continuous function on a compact set, then existence of a minimizer would be easy by the Weierstrass theorem.

Here the situation is more complicated:
\begin{itemize}
    \item the space is not finite-dimensional;
    \item the topology is usually convergence in probability, i.e. convergence in $L^0$;
    \item compactness is replaced by a weaker sequential compactness principle;
    \item the function
    \[
        h\mapsto \E[J(h)]
    \]
    is convex and lower semicontinuous.
\end{itemize}

Still, the strategy is morally the same.

\subsubsection*{Step 1: Take a minimizing sequence}

Choose a sequence
\[
    (h_n)_{n\in\mathbb{N}}\subseteq \D(z)
\]
such that
\[
    \E[J(h_n)]\downarrow j(z).
\]
This means that the sequence gets closer and closer to the optimal value.

\subsubsection*{Step 2: Use convex compactness}

In convex optimization, ordinary subsequences are sometimes not enough. Instead, we often use convex combinations.

So we try to find convex combinations
\[
    \tilde h_n
    \in
    \operatorname{conv}(h_n,h_{n+1},h_{n+2},\dots)
\]
such that
\[
    \tilde h_n \to h_z^*
\]
in probability.

The limit $h_z^*$ is the candidate optimizer.

\begin{opomba}
This is why convexity matters. Even if the original sequence does not converge, suitable convex combinations often do.
\end{opomba}

\subsubsection*{Step 3: Show the limit stays in the domain}

We need to show
\[
    h_z^*\in\D(z).
\]
This usually follows from the fact that $\D(z)$ is convex and closed in the appropriate topology.

\subsubsection*{Step 4: Use lower semicontinuity}

Since
\[
    h\mapsto \E[J(h)]
\]
is lower semicontinuous, we get
\[
    \E[J(h_z^*)]
    \leq
    \liminf_{n\to\infty}\E[J(\tilde h_n)].
\]
Because the $\tilde h_n$ are made from a minimizing sequence, the right-hand side is equal to $j(z)$. Hence
\[
    \E[J(h_z^*)]\leq j(z).
\]
But $j(z)$ is the infimum, so automatically
\[
    \E[J(h_z^*)]\geq j(z).
\]
Therefore
\[
    \E[J(h_z^*)]=j(z).
\]
So $h_z^*$ is a minimizer.

\subsubsection*{Step 5: Uniqueness}

The function $J$ is strictly convex. Therefore, if there were two different minimizers
\[
    h_1,h_2\in\D(z),
\]
then their average
\[
    \frac{h_1+h_2}{2}
\]
would also belong to $\D(z)$ by convexity. Strict convexity would imply
\[
    J\left(\frac{h_1+h_2}{2}\right)
    <
    \frac{1}{2}J(h_1)+\frac{1}{2}J(h_2)
\]
on the set where $h_1\neq h_2$.

Taking expectations would give a strictly smaller value than the minimum, which is impossible. Hence the minimizer is unique.

\begin{opomba}
So the slogan is:

\[
    \text{existence comes from compactness/closedness,}
\]
and
\[
    \text{uniqueness comes from strict convexity.}
\]
\end{opomba}

\subsection{Summary of the symbols}

\begin{itemize}
    \item $x$: initial capital.
    \item $\theta$: trading strategy.
    \item $G(\theta)$: gains from trading.
    \item $V(x,\theta)=x+G(\theta)$: wealth process.
    \item $\V(x)$: all admissible nonnegative wealth processes starting from $x$.
    \item $U$: utility function.
    \item $u(x)$: maximal expected utility starting from $x$.
    \item $\C(x)$: all nonnegative terminal payoffs superreplicable from $x$.
    \item $J$: convex conjugate / Legendre transform of $U$.
    \item $z$: initial value of the dual variable.
    \item $Z$: supermartingale deflator.
    \item $\Z(z)$: all supermartingale deflators starting from $z$.
    \item $j(z)$: dual value function.
    \item $\D(z)$: abstract terminal dual domain.
    \item $h_z^*$: optimizer of the dual problem.
\end{itemize}

\subsection{Big picture}

The whole theory can be summarized as:

\[
    \boxed{
    u(x)
    =
    \sup_{f\in\C(x)}
    \E[U(f)]
    }
\]
is the primal problem, and

\[
    \boxed{
    j(z)
    =
    \inf_{h\in\D(z)}
    \E[J(h)]
    }
\]
is the dual problem.

The bridge between them is the inequality
\[
    U(a)\leq J(y)+ay.
\]
Applied to
\[
    a=V_T,
    \qquad
    y=Z_T,
\]
this gives
\[
    \E[U(V_T)]
    \leq
    \E[J(Z_T)]
    +
    \E[Z_TV_T].
\]
The supermartingale property gives
\[
    \E[Z_TV_T]\leq xz.
\]
Therefore
\[
    u(x)\leq j(z)+xz.
\]

This is the basic weak duality inequality. Under additional assumptions, one can usually prove equality in the appropriate conjugate sense:
\[
    u(x)
    =
    \inf_{z>0}\bigl(j(z)+xz\bigr).
\]
